\section{Modules développés}

\begin{figure}[H]
    \centering
    \includegraphics[width=12cm]{diagram/class_diagram.png}
    \caption{Diagramme de classe du projet \newline \textit{N'include que les module développé par nos soins}}
    \label{fig:class_diagram}
\end{figure}

\subsection{Module capteur ultrasonique (ULTRA\_SONIC)}
Ce module implémente un pattern Singleton pour gérer le capteur HC-SR04. Il offre les fonctionnalités suivantes:

\begin{lstlisting}[caption={Module capteur ultrasonique}, label={lst:ultrasonicmodule}]
void ULTRA_SONIC_Init(void)
{
    // Initialisation du capteur ultrasonique
    HAL_TIM_PWM_Start(&htim4, TIM_CHANNEL_2);
    GPIO_RegisterGPIO_EXTICallback(gpio_trigger_Callback, SonicSensor_Echo_Pin);
}

float ULTRA_SONIC_GetDistance(void)
{
    return distance; 
}

void ULTRA_SONIC_test(void){
    float distance = ULTRA_SONIC_GetDistance();
    char buffer[50];
    int len = snprintf(buffer, sizeof(buffer), "Distance: %f \r\n", distance);
    HAL_UART_Transmit(&huart2, (uint8_t *)buffer, len, HAL_MAX_DELAY);
    HAL_Delay(1000);
}
\end{lstlisting}

Le principe de fonctionnement est le suivant:
\begin{enumerate}
    \item Envoi d'une impulsion de 10 µs sur le pin TRIG
    \item Attente du signal ECHO et mesure de sa durée via interruption
    \item Calcul de la distance selon la formule: $distance = \frac{duree \times vitesse\_son}{2}$
\end{enumerate}

\subsection{Module servomoteur (SERVO)}
Ce module implémente également un pattern Singleton pour gérer le servomoteur:

\begin{lstlisting}[caption={Module servomoteur}, label={lst:servomodule}]
#define MIN_DUTY 2000
#define MAX_DUTY 4000

void SERVO_Init(void)
{
    // Initialisation du servomoteur
    HAL_TIM_PWM_Start(&htim3, TIM_CHANNEL_2);
    __HAL_TIM_SetCompare(&htim3, TIM_CHANNEL_2, MIN_DUTY);
}

void SERVO_set_servo_percentage(uint8_t percentage)
{
    // Configuration du rapport cyclique selon le pourcentage
    uint32_t duty_cycle = MIN_DUTY + (MAX_DUTY - MIN_DUTY) * percentage / 100;
    __HAL_TIM_SetCompare(&htim3, TIM_CHANNEL_2, duty_cycle);
}
\end{lstlisting}

Le servomoteur est contrôlé par un signal PWM dont la largeur d'impulsion détermine la position angulaire:
\begin{itemize}
    \item MIN\_DUTY (2000) correspond à la position 0°
    \item MAX\_DUTY (4000) correspond à la position 180°
\end{itemize}

\subsection{Module LED (LEDS)}
Un module simple pour gérer les LEDs d'indication de mode:

\begin{lstlisting}[caption={Interface du module LED}, label={lst:ledmodule}]
void LEDS_Init(void);
void LEDS_SetMode(uint8_t mode);
\end{lstlisting}

\subsection{Module communication série}
Le module de communication série permet d'envoyer la distance mesurée et de recevoir les commandes:

\begin{lstlisting}[caption={Exemple de communication série}, label={lst:serialcomm}]
// Envoi de la distance
float distance = ULTRA_SONIC_GetDistance();
char buffer[50];
int len = snprintf(buffer, sizeof(buffer), "Distance: %f \r\n", distance);
HAL_UART_Transmit(&huart2, (uint8_t *)buffer, len, HAL_MAX_DELAY);
\end{lstlisting}

\subsection{Module bouton poussoir (BUTTON)}
Le module de gestion du bouton poussoir permet de changer le mode de fonctionnement du système. Il utilise une interruption externe pour détecter les pressions sur le bouton:

Le module BUTTON gère un bouton poussoir de manière autonome en utilisant une interruption externe. Il se base sur deux variables globales principales : une variable indiquant si le bouton a été pressé et une autre mémorisant le dernier instant d’activation, ce qui est crucial pour le contrôle du rebond.

\begin{enumerate}
    \item Gestion du rebond
    \begin{itemize}
        \item Lorsqu'une pression sur le bouton est détectée, une interruption est générée.
        \item La fonction de rappel vérifie si l'intervalle écoulé depuis le dernier appui dépasse un seuil (200 ms).
        \item Si le délai n'est pas respecté, l'appui est ignoré pour éviter les faux signaux.
    \end{itemize}
    \item Mise à jour de l’état et communication
    \begin{itemize}
        \item En cas d’appui valide, la fonction de rappel met à jour la variable d’état (passant à 1).
        \item Elle enregistre l’instant courant via un appel système (par exemple, HAL\_GetTick()).
        \item Un message est envoyé via l’interface UART pour notifier le reste du système de l’événement.
    \end{itemize}
    \item Initialisation et configuration
    \begin{itemize}
        \item Avant d’être opérationnel, le module doit être initialisé.
        \item Une fonction dédiée configure le bouton pour qu’il génère une interruption à chaque pression.
        \item Elle associe le pin correspondant au bouton à la fonction de rappel, garantissant ainsi que l’appui sur le bouton déclenche automatiquement la gestion d’événement.
    \end{itemize}
    \item Lecture et réinitialisation de l’état
    \begin{itemize}
        \item Une fonction de lecture permet à d'autres parties du programme de vérifier si le bouton a été pressé.
        \item Une fois l’événement traité, une fonction de réinitialisation remet la variable d’état à zéro pour préparer la détection de futures pressions.
    \end{itemize}
\end{enumerate}
